
\section{Spin and orbital angular-momentum densities}

The intrinsic angular-momentum density is obtained by taking moments
of the linear-momentum density about the energy-flux centroid.  On the
detector screen, let ${\bf r}_\perp=(x_1,x_2,0)$ and
${\bf R}_c(\omega)=(X_c(\omega),Y_c(\omega),0)$.  Its third component in the
one-sided spectrum is
\begin{align}
  \frac{d\mathcal L^{\mathrm{int}}_{3,+}}{d\omega}
  &=4\pi\epsilon_0\left[
    ({\bf r}_\perp-{\bf R}_c)\times
    \operatorname{Re}\bigl({\bf E}(\omega,{\bf x})
    \times{\bf B}(\omega,{\bf x})^*\bigr)
  \right]_3,
  \qquad \omega>0.
  \label{eq:one-sided-intrinsic-kinetic-angular-momentum-density}
\end{align}

We adopt the dual-symmetric spin--orbital separation for a source-free
electromagnetic field \cite{bliokh_conservation_2014}.  For the one-sided
Fourier spectrum, the spin density of the third component is
\begin{align}
  \frac{d\mathcal L^{\mathrm{spin}}_{3,+}}{d\omega}
  =\frac{2\pi\epsilon_0}{\omega}\operatorname{Im}\left[
    \left({\bf E}^*\times{\bf E}\right)_3
    +c^2\left({\bf B}^*\times{\bf B}\right)_3
  \right],
  \qquad \omega>0.
  \label{eq:one-sided-spin-angular-momentum-density}
\end{align}
Here and below, all fields in a spectral expression are evaluated at
$(\omega,{\bf x})$.  One local representative of the orbital density is the
residual obtained by subtracting the spin density from the intrinsic kinetic
density:

\begin{align}
  \frac{d\mathcal L^{\mathrm{orb,res}}_{3,+}}{d\omega}
  =\frac{d\mathcal L^{\mathrm{int}}_{3,+}}{d\omega}
  -\frac{d\mathcal L^{\mathrm{spin}}_{3,+}}{d\omega}.
  \label{eq:one-sided-intrinsic-orbital-angular-momentum-density-dif}
\end{align}
The dual-symmetric canonical orbital density is instead expressed through the
generator of rotations about $Ox_3$.  On the detector screen, introduce polar
coordinates about the energy-flux centroid at each frequency by
\begin{align}
  x_1-X_c=\rho\cos\phi,
  \qquad
  x_2-Y_c=\rho\sin\phi.
  \label{eq:screen-polar-coordinates-about-centroid}
\end{align}
The generator of rotations about the third axis is then
\begin{align}
  \mathscr L_3
  =(x_1-X_c)\partial_2-(x_2-Y_c)\partial_1
  =\frac{\partial}{\partial\phi}.
  \label{eq:intrinsic-orbital-rotation-generator}
\end{align}
Consequently, the canonical intrinsic orbital density on the screen is
\begin{align}
  \frac{d\mathcal L^{\mathrm{orb,can}}_{3,+}}{d\omega}
  =\frac{2\pi\epsilon_0}{\omega}\operatorname{Im}\left[
    \sum_{a=1}^{3}E_a^*\mathscr L_3E_a
    +c^2\sum_{a=1}^{3}B_a^*\mathscr L_3B_a
  \right],
  \qquad \omega>0,
  \label{eq:one-sided-intrinsic-orbital-angular-momentum-density}
\end{align}
or, equivalently,
\begin{align}
  \frac{d\mathcal L^{\mathrm{orb,can}}_{3,+}}{d\omega}
  =\frac{2\pi\epsilon_0}{\omega}\operatorname{Im}\left[
    \sum_{a=1}^{3}E_a^*\frac{\partial E_a}{\partial\phi}
    +c^2\sum_{a=1}^{3}B_a^*\frac{\partial B_a}{\partial\phi}
  \right],
  \qquad \omega>0.
  \label{eq:one-sided-intrinsic-orbital-angular-momentum-density-azimuthal}
\end{align}

For clarity, the four quantities needed in the numerical calculation can all
be expressed directly through $F^{\alpha\beta}(\omega,{\bf x})$.  First, the
intrinsic  density is
\begin{highlightbox}[title={Intrinsic angular-momentum density from the Faraday tensor}]
  \begin{align}
    \frac{d\mathcal L^{\mathrm{int}}_{3,+}}{d\omega}
    =-4\pi\epsilon_0c\left\{
      (x_1-X_c)\operatorname{Re}\left[
      \sum_{j=1}^{3}F^{j0}(F^{2j})^*\right]-(x_2-Y_c)\operatorname{Re}\left[
      \sum_{j=1}^{3}F^{j0}(F^{1j})^*\right]
    \right\}.
    \label{eq:one-sided-intrinsic-kinetic-angular-momentum-density-Faraday}
  \end{align}
\end{highlightbox}
\begin{highlightbox}[title={Spin angular-momentum density from the Faraday tensor}]
  \begin{align}
    \frac{d\mathcal L^{\mathrm{spin}}_{3,+}}{d\omega}
    =\frac{2\pi\epsilon_0c^2}{\omega}\operatorname{Im}\left[(F^{10})^*F^{20}-(F^{20})^*F^{10}+(F^{13})^*F^{23}-(F^{23})^*F^{13}
    \right].
    \label{eq:one-sided-spin-angular-momentum-density-Faraday}
  \end{align}
\end{highlightbox}
\begin{highlightbox}[title={Orbital angular-momentum densities from the Faraday tensor}]
  The orbital density obtained as the residual of the intrinsic kinetic
  density is
  \begin{align}
    \frac{d\mathcal L^{\mathrm{orb,res}}_{3,+}}{d\omega}
    =\frac{d\mathcal L^{\mathrm{int}}_{3,+}}{d\omega}
    -\frac{d\mathcal L^{\mathrm{spin}}_{3,+}}{d\omega}.
    \label{eq:one-sided-residual-orbital-angular-momentum-density-Faraday}
  \end{align}
  In the dual-symmetric canonical separation, the orbital density is instead
  evaluated directly from the azimuthal variation of the field,
  \begin{align}
    \frac{d\mathcal L^{\mathrm{orb,can}}_{3,+}}{d\omega}
    =\frac{2\pi\epsilon_0c^2}{\omega}\operatorname{Im}\left[
      \sum_{a=1}^{3}(F^{a0})^*\frac{\partial F^{a0}}{\partial\phi}
      +\sum_{1\leq a<b\leq3}(F^{ab})^*
      \frac{\partial F^{ab}}{\partial\phi}
    \right].
    \label{eq:one-sided-canonical-orbital-angular-momentum-density-Faraday}
  \end{align}
  Here $\omega>0$, every Faraday component is evaluated at
  $(\omega,{\bf x})$, and $\partial_\phi$ is taken at fixed $\rho$ about the
  energy-flux centroid.
\end{highlightbox}

The sum of the spin and canonical orbital densities defines the canonical
intrinsic angular-momentum density,
\begin{align}
  \frac{d\mathcal L^{\mathrm{int,can}}_{3,+}}{d\omega}
  =\frac{d\mathcal L^{\mathrm{spin}}_{3,+}}{d\omega}
  +\frac{d\mathcal L^{\mathrm{orb,can}}_{3,+}}{d\omega}.
  \label{eq:one-sided-intrinsic-angular-momentum-density-split}
\end{align}
The residual and canonical orbital densities need not coincide point by point:
the kinetic and canonical momentum densities differ by a spatial-divergence
term.  Their integrals over the complete screen agree when the associated
boundary term vanishes.  No relation between ${\bf E}$ and ${\bf B}$ beyond
Maxwell's equations has been imposed.
